% A.61 candidate. Unregistered pending full independent audits.
\subsection{An unbounded Heisenberg family from a divided differential operator}

\begin{statement}{A.61}{Proved}{Both assertions for every inert-prime Heisenberg cover}
For every prime \(p\equiv3\pmod4\), let \(C_p\) be the
smooth proper normalization of the function field
\[
 E:y^2=x^3-x,\qquad z^p=\psi_p(x),
 \qquad q:C_p\longrightarrow E,
                                                        \tag{A.61.1}
\]
where \(\psi_p\) is the usual division polynomial. Then
\[
 Y_p=C_p\setminus q^{-1}(E[p])
           \xrightarrow{[p]q}U
                                                        \tag{A.61.2}
\]
is an actual degree-\(p^3\) Heisenberg cover, with
\(p^2\) cusps and proper genus
\(1+p^2(p-1)/2\). It satisfies both assertions of
Task A.5 at every coefficient prime: all nontrivial
arithmetically normalized procyclic subgroups are
cuspidal, and every permitted cusp packet and indicated
open derived subgroup has full fixed group exactly
the selected cusp inertia.

At a place above \(p\), an explicit stable model has
one supersingular elliptic main component and \(p\)
positive-genus tails, each geometrically isomorphic to
\[
                         V^p-V=U^{p+1},
 \qquad g=\frac{p(p-1)}2.                              \tag{A.61.3}
\]
Every other component is rational, the dual graph is
a tree, and all \(p^2\) cusps lie on a rational marked
tree over the main point above \(O\). The proof is
uniform in \(p\); no finite polynomial pattern is
being substituted for it.
\end{statement}

\begin{proof}
We establish a uniform local coefficient identity before
constructing the model. Throughout its proof all reduced
functions and constants are defined over \(\F_p\),
even when their zeros are subsequently separated over
a finite extension.

\paragraph{Frobenius, and the integral divided polynomial.}
Write \(f=x^3-x\). Pairing \(x\) with \(-x\) in the
quadratic-character point count gives
\(\#E(\F_p)=p+1\), since \(-1\) is a nonsquare.
The Frobenius endomorphism \(F\) therefore satisfies
\[
                           F^2=[-p].                  \tag{A.61.4}
\]
This is the actual endomorphism identity of
\cite[Theorem V.2.3.1(b), p.142]{fghSilvermanAEC}, not
only a characteristic-polynomial statement on one
Tate module. It also proves supersingularity: the
multiplication-by-\(p\) morphism is purely inseparable
of degree \(p^2\), so its geometric kernel is \(\{O\}\).

The exact division formulas are
\[
 x\circ[p]=\frac{\phi_p}{\psi_p^2},\qquad
 \phi_p=x\psi_p^2-\psi_{p+1}\psi_{p-1}.
                                                        \tag{A.61.5}
\]
The polynomial \(\phi_p\) is monic of degree \(p^2\);
the numerator and denominator are relatively prime on
a nonsingular Weierstrass curve. These assertions and
the division recurrence are
\cite[Exercise III.3.7(a)--(d), pp.105--106]{fghSilvermanAEC}.
Reducing (A.61.5) and using (A.61.4) gives
\(\overline\phi_p=x^{p^2}\overline\psi_p^{\,2}\).
The denominator is nonzero: otherwise relative primality
would fail for the nonconstant monic numerator. Comparing
degrees and leading coefficients shows
\(\overline\psi_p=c\), a nonzero constant with \(c^2=1\).

In fact \(\psi_n(0)=(-1)^{(n-1)/2}\) for odd \(n\).
This follows directly by induction from the odd division
recurrence at the two-torsion point \((0,0)\): the even
division functions vanish there; for an odd middle index
only the first product remains, and for an even middle
index only the negative second product remains. The
initial values are \(\psi_1(0)=1\), \(\psi_3(0)=-1\).
Thus in the present case \(c=-1\), and
\[
             h=\frac{\psi_p+1}{p}\in\Z[x],
       \qquad \deg h=\frac{p^2-1}{2},
       \qquad \text{leading coefficient }1.           \tag{A.61.6}
\]
We retain the notation \(c=-1\) in the identities below
to make their origins visible.

\paragraph{The invariant derivation and its divided power.}
Let
\[
 \mathcal A=\Z_{(p)}[x,y]/(y^2-x^3+x),\qquad
 \omega=\frac{dx}{2y},\qquad
 D(x)=2y,\quad D(y)=3x^2-1.                            \tag{A.61.7}
\]
On the smooth elliptic curve the derivation \(D\) is
the nowhere-vanishing invariant vector field dual to
\(\omega\). Its reduction satisfies \(D^p=0\).
Here is a direct verification for this curve. In
characteristic \(p\), the \(p\)-th power of a derivation
is a derivation; it is still translation-invariant,
so it is a scalar multiple of \(D\). Over \(\F_{p^2}\),
the automorphism \(\rho(x,y)=(-x,iy)\) conjugates
\(D\) to \(-iD\), and therefore \(D^p\) to
\((-i)^pD^p=iD^p\). A scalar multiple of \(D\)
instead transforms by \(-i\). Since \(p\) is odd,
these are different, proving \(D^p=0\).

Consequently \(D^p\mathcal A\subset p\mathcal A\),
and the operation
\[
                \mathcal E(\overline b)
                  =\overline{D^p b/p}
                  \quad(b\in\mathcal A)               \tag{A.61.8}
\]
is a well-defined \(\F_p\)-linear operator on
\(\overline{\mathcal A}\). Changing the lift by
\(p b_1\) changes the right side by
\(\overline{D^p b_1}=0\). It commutes with \(D\).
It is not being treated as a derivation. Its essential
identity is
\[
                    \mathcal E(b^p)=-(Db)^p.           \tag{A.61.9}
\]
To prove it, expand \(D^p\) on a product of \(p\)
identical factors. The distribution with one derivative
on every factor contributes \(p!(Db)^p\), which after
division by \(p\) reduces to \(-(Db)^p\). The
distributions with all \(p\) derivatives on one factor
sum to \(p b^{p-1}D^p b\), and vanish after division
and reduction. Every other distribution has all entries
less than \(p\), so its multinomial coefficient is
divisible by \(p\). Its orbit under cyclic permutation
of the \(p\) identical factors has length \(p\);
the additional orbit factor makes it divisible by
\(p^2\). These are all cases, proving (A.61.9).

If a reduced function \(R\) satisfies \(DR=a\in\F_p\),
the same expansion on \(p+1\) factors gives
\[
 \mathcal E(R^{p+1})
       =R^p\mathcal E(R)-R a^p.                       \tag{A.61.10}
\]
Indeed the single \(p\)-derivative terms give the first
term. The distributions with \(p\) first derivatives
give the second by Wilson's congruence. Every other
term contains some \(D^jR\), \(2\le j<p\), which
vanishes modulo \(p\), in addition to the factor
\(p\) already present in the multinomial coefficient.
This proves (A.61.10) with arbitrary integral lifts.

\paragraph{The division-polynomial differential identity.}
In characteristic zero one has the exact rational-function
identity
\[
 \psi_p D^2\psi_p-(D\psi_p)^2
       =p^2\psi_{p+1}\psi_{p-1}.                      \tag{A.61.11}
\]
For completeness, equivalently
\(D^2\log\psi_p=p^2(x-x\circ[p])\).
In a formal coordinate \(s\) with \(ds=\omega\),
the principal part on either side at every nonzero
\(p\)-torsion point is \(-s^{-2}\); at \(O\) it is
\((p^2-1)s^{-2}\). There are no other poles. Their
difference is therefore a constant. At \(O\),
\(x=s^{-2}+O(s^2)\). Moreover \(\psi_p\) is an
even polynomial in \(x\): pullback by \([i]\)
preserves its divisor and its leading coefficient,
since \((p^2-1)/2\) is even. Thus its next-to-leading
coefficient vanishes, and
\(\psi_p=p s^{1-p^2}(1+O(s^4))\).
Both constant terms in the displayed logarithmic
identity are zero. This proves it, and (A.61.5)
gives (A.61.11). The formal Laurent assertion for
\(x\) also follows immediately from
\((Dx)^2=4(x^3-x)\); no transcendental comparison
is required.

Substitute \(\psi_p=-1+ph\) in (A.61.11).
Reduction first proves \(D^2\overline h=0\).
Hence
\[
                 D\overline h=R^p
               \quad\text{for a unique }R\in\F_p(E).
                                                        \tag{A.61.12}
\]
The kernel of a nonzero derivation on the one-variable
function field \(\F_p(E)\) is its subfield of
\(p\)-th powers. The function on the left is regular
away from \(O\), so normality makes \(R\) regular
there too. Its pole has order exactly \(p\): the
leading coefficient and degree in (A.61.6) give
\(D\overline h\) a pole of order \(p^2\).
The function \(R\) is odd under \([-1]\), since
\(\overline h\) is a polynomial in \(x\).

Put \(T=D^2h/p\in\mathcal A\). Its integrality
also follows from (A.61.11), since the reduction of
\(\psi_p\) is a unit. Reducing the full identity and
using (A.61.4)--(A.61.5) yields
\[
 c\overline T-R^{2p}=c^2(x-x^{p^2}),
                         \qquad c=-1.                 \tag{A.61.13}
\]
Apply \(D^{p-2}\), which is a positive power even
when \(p=3\). The \(p\)-th-power terms disappear,
and (A.61.8) gives
\[
                   \mathcal E(\overline h)
                           =cD^{p-2}x.                \tag{A.61.14}
\]
Apply \(D\) and use (A.61.9) and commutation:
\[
                      -(DR)^p=cD^{p-1}x.               \tag{A.61.15}
\]
The function \(D^{p-1}x\) is killed by \(D\) and
has pole order at most \(p+1<2p\). Its \(p\)-th
root can have pole order at most one, so it is
constant on the genus-one curve. Thus
\[
                         DR=a\in\F_p^\times.          \tag{A.61.16}
\]
The constant is nonzero: otherwise \(R\) would be
a \(p\)-th power whose root has a single simple
pole, which is impossible on an elliptic curve.
Because \(a\in\F_p\), equation (A.61.15) says
\(cD^{p-1}x=-a\). Therefore
\(cD^{p-2}x+R\) is killed by \(D\) and has pole
order at most \(p\). Its \(p\)-th root is again
constant. Both summands are odd under \([-1]\),
so that constant is zero. We have proved
\[
               R=-cD^{p-2}x,
                \qquad \mathcal E(\overline h)=-R.    \tag{A.61.17}
\]
This identifies the actual divided polynomial's
derivative; \(D^2\overline h=0\) alone would not
have sufficed.

\paragraph{The hidden degree-p Taylor coefficient vanishes.}
By (A.61.12) and (A.61.16), there is a unique
\(S\in\F_p(E)\) such that
\[
            \overline h=\frac{R^{p+1}}a+S^p.           \tag{A.61.18}
\]
Its only possible pole is at \(O\), of order at
most \(p+1\). Since \(\mathcal E\) commutes with
\(D\) and \(DR=a\), the function
\(\mathcal E(R)\) is killed by \(D\). Its pole
order is at most \(2p\): take an integral lift of
\(R\) of pole order at most \(p\), and apply
\(D^p/p\). Consequently
\[
           \mathcal E(R)=A^p,
                      \qquad A\in L(2O)(\F_p).        \tag{A.61.19}
\]
Apply \(\mathcal E\) to (A.61.18), using
(A.61.9)--(A.61.10), (A.61.17), and
\(a^{p-1}=1\). It gives
\[
 -R=\left(\frac{RA}{a}\right)^p-R-(DS)^p,
                  \qquad DS=\frac{RA}{a}.             \tag{A.61.20}
\]
The last equality follows from injectivity of
Frobenius. Every constant and every function in
(A.61.12)--(A.61.20) is defined over \(\F_p\);
in particular the use of \(a^{p-1}=1\) is justified.

The differential \(dR=a\omega\) never vanishes.
Since \(R\) has one pole of order \(p\), it has
exactly \(p\) distinct simple zeros, all away from
\(O\). At any such point \(P\), equation
(A.61.20) says \(dS(P)=0\). For a local parameter
\(t\), equation (A.61.18) now gives
\[
 \overline h(P+t)-\overline h(P)
       =\frac{R(P+t)^{p+1}}a
          +(S(P+t)-S(P))^p,
 \qquad \operatorname{ord}_P(\overline h-\overline h(P))=p+1.
                                                        \tag{A.61.21}
\]
Indeed the first term has exactly that order and
the second has order at least \(2p>p+1\).
Thus every Taylor coefficient of degrees one
through \(p\), including the degree invisible
to ordinary differentiation, is zero.

More explicitly \(R=yQ(x)\), with
\(\deg Q=(p-3)/2\). Equation (A.61.16) is
\[
       (3x^2-1)Q+2(x^3-x)Q'=a.                         \tag{A.61.22}
\]
It makes \(Q\) squarefree and disjoint from
\(x(x^2-1)\). The zeros of \(R\) consist of
the three finite two-torsion points and
\(p-3\) other points, two over each root of
\(Q\). There are no omitted critical centres.

\paragraph{Actual algebraic main and tail charts.}
Choose \(\pi^p=p\) and set \(z=-1+\pi u\).
The exact expansion of (A.61.1), after division
by \(p\), has reduction \(u^p=\overline h\).
Its differential is nonzero by (A.61.12), so its
normalized special function field is
\(\F_p(E)^{1/p}\). It gives a genus-one main
component mapping radicially with degree \(p\)
to the good elliptic special fibre. Its special
generic point is smooth, using the nonzero
differential and the elliptic equation.

Choose the \(p\) residue centres just found over
a finite \emph{unramified} extension \(R_0/\Z_p\).
For a non-two-torsion centre, lift its \(x\)-coordinate
as a simple root of an integral lift of \(Q\), then
lift \(y\) by the unit-derivative elliptic equation.
At a two-torsion point take \((a,0)\) with
\(a=0,1,-1\). Equation (A.61.21) therefore implies
divisibility of every lower integral Taylor
coefficient by \(p\), not merely by a ramified
uniformizer.

Take \(\eta^{p-1}=-p\), \(\rho^{p+1}=\eta\),
\(z_0^p=\psi_p(a)\), and put
\(z=z_0(1+\eta V)\). The exact equation is
\[
 \begin{aligned}
 \mathcal B_\eta(V)
    &=\sum_{j=1}^{p-1}\frac{\binom pj}{p}
                       \eta^{j-1}V^j-V^p,\\
 z_0^p\mathcal B_\eta(V)
    &=\frac{h(x)-h(a)}{\eta}.                          \tag{A.61.23}
 \end{aligned}
\]
At a centre \((a,y_0)\) with \(y_0\ne0\), set
\(x=a+\rho U\), \(y=y_0+\rho W\). The elliptic
equation becomes
\[
 2y_0W+\rho W^2
             =(3a^2-1)U+3a\rho U^2+\rho^2U^3.         \tag{A.61.24}
\]
The leading term on the right of (A.61.23) is
a nonzero multiple of \(U^{p+1}\).
For a lower degree \(1\le j<p+1\), its valuation
after division by \(\eta\) is at least
\[
                   1+\frac{j-(p+1)}{p^2-1}>0;
                                                        \tag{A.61.25}
\]
for a higher degree it is at least
\((j-(p+1))/(p^2-1)>0\).

At a two-torsion centre use instead
\(y=\rho U\), \(x=a+\rho^2W\), with
\(\delta=3a^2-1\) a unit. Then
\[
                 U^2=\delta W+3a\rho^2W^2+\rho^4W^3.
                                                        \tag{A.61.26}
\]
The first nonzero degree of
\(\overline h(a+X)-\overline h(a)\) is
\((p+1)/2\), since \(y\) is the local parameter.
The same valuation calculation, replacing \(j\)
by \(2j\), leaves a nonzero multiple of
\(U^{p+1}\). All other terms have strictly
positive valuation.

The special equations are therefore
\(V^p-V=c_PU^{p+1}\), \(c_P\ne0\).
The partials in \(W,V\) are units. These are
actual smooth algebraic charts with integral
special curves, and their generic coordinate
changes are invertible on (A.61.1). The
Artin--Schreier polynomial is irreducible because
its pole order \(p+1\) is prime to \(p\).
Over the algebraic residue field rescaling
\(U\) gives (A.61.3).

At its unique point over infinity put \(t=U^{-1}\).
One has \(v(t)=p\), \(v(V)=-(p+1)\), and
\(t^pV^{p-1}\) is a uniformizer. Its difference
under \(V\mapsto V+b\), \(b\ne0\), has
valuation \(p+2\). The lower jump is \(p+1\),
the different is \((p-1)(p+2)\), and Hurwitz
gives genus \(p(p-1)/2\). Its \(p\)-rank is
zero by the one-branch-point Deuring--Shafarevich
formula, \cite[Corollary 1.8]{fghCrew1984}.

\paragraph{The actual cover, all components, and every cusp.}
The division divisor is
\(D_{E[p]}-p^2(O)\). Thus \(q\) is connected
and totally ramified at every point of \(E[p]\),
with orders one and \(1-p^2\), both prime to
\(p\). It has genus \(1+p^2(p-1)/2\), and
(A.61.2) has degree \(p^3\) and \(p^2\) cusps.
Every translation by \(T\in E[p]\) lifts:
the divisor of its division-function ratio is
\(p^2(O)-p^2(-T)\), which is \(p\) times
the principal divisor \(p(O)-p(-T)\).
Adjoining roots of the finitely many remaining
constants gives the actual lifts over a number
field. The deck group is a central extension
of \(E[p]\) by \(\mu_p\). It is nonabelian
because peripheral commutator inertia has
order \(p\). The automorphism
\((x,y,z)\mapsto(-x,iy,z)\) fixes \(\mu_p\)
and acts by \([i]\) on \(E[p]\).
The \(p\)-th-power map \(E[p]\to\mu_p\)
is additive, since \(p\mid\binom p2\), and
is \([i]\)-invariant. As \([i]-1\) is
invertible on \(E[p]\), this map is zero.
The group is therefore the exponent-\(p\)
Heisenberg group of order \(p^3\), exactly as
in the individual verification of A.60.

The positive-genus valuations furnished by
the main chart and the \(p\) tail charts
persist in the stable model. They are distinct:
the tail restrictions have the \(p\) different
finite elliptic centres, and the main restriction
is the generic point. The precise common-regular-
resolution and positive-genus persistence argument
is that of A.52, A.58 and A.60, using
\cite[Theorem 9.2.2, Remark 9.2.6 and Theorem 9.3.21]
{fghLiu}. Their genera already sum to
\[
                     1+p\frac{p(p-1)}2
                        =1+\frac{p^2(p-1)}2.          \tag{A.61.27}
\]
The semistable genus formula forces compact
type and makes every remaining component
rational. No unramified-cover theorem is
being applied to the ramified map \(q\).

The actual map to the good elliptic model
extends by resolving its indeterminacy and
contracting rational exceptional components,
as in A.58. Its main component already has
degree \(p\). Constancy of the degree of a
pulled-back line bundle forces every other
component to map to a point. Every off-main
subtree has one attachment point and maps
constantly to that image.

All generic \(p\)-torsion points specialize
to \(O\), by (A.61.4). Each has a unique
ramified cusp above it. The positive-genus
tails lie above zeros of \(R\), all different
from \(O\), and hence no cusp lies on those
tails or their connecting subtrees. The
radicial main map has one point over \(O\).
Separating the \(p^2\) sections therefore
puts every mark on a rational marked tree
attached there. No mark remains at its
attachment node. Every chosen cusp has a
rational joining path to another, including
both endpoints. Statements A.41--A.42 now
give the full fixed-inertia conclusion at
every coefficient prime and every permitted
packet. Statement A.59 gives every normalizer
because the actual geometric Galois group
is a \(p\)-group.

The accompanying exact program checks the
division identities, divided operators,
Frobenius roots \(R,A,S\), the identity
\(DS=RA/a\), and all chart coefficients at
several inert primes. It is a certificate
for the formulas; the uniform argument
(A.61.8)--(A.61.21) is what proves the
assertion for every inert prime.
\end{proof}
